\section{Introduction}\label{sec:intro}

\subsection{Motivation}
An industrial plant is a poor fit for the assumptions behind most identity systems. Its devices live for fifteen or twenty years, speak protocols that predate public-key cryptography, sit behind process-critical networks that cannot be rebooted at will, and are administered by several organizations at once: the operator who runs the line, the original manufacturer who ships firmware, an integrator who commissioned the cell, and sometimes an outside maintenance contractor. Guidance for operational technology (OT) security is explicit that availability and safety, not confidentiality, usually dominate the design space, and that a control that interrupts a process can be worse than the attack it prevents \cite{stouffer2023sp80082}. At the same time the attack surface has widened, because IT--OT convergence connects field assets to analytics, remote support and cloud services. Zero Trust, which drops trust based on network location and demands that every access be authenticated, authorized and re-evaluated \cite{nist800207,fernandez2024critical,phiayura2023comprehensive}, is a natural response, and early proposals adapt it to critical OT systems \cite{song2025ztot,zanasi2024flexible}.

Identity is where the difficulty concentrates. A conventional public-key infrastructure gives each device a certificate and relies on revocation lists or short certificate lifetimes to withdraw it. In a multi-party plant, however, there is rarely one party that every other party is willing to let revoke a controller, and the delay between the decision to withdraw a device and the moment every gateway honours it is exactly the window in which a compromised or recalled device keeps working. Blockchain-based identity management addresses the first problem by anchoring identifiers, keys and status on a tamper-evident ledger that several organizations operate together, following the decentralized-identifier (DID) model \cite{w3cdid2022,hosseini2023blockchain,aguru2025ssi,folino2025puf}. Recent industrial work builds authentication on this foundation \cite{zhai2025chebyshev}, and surveys of blockchain for Industrial IoT (IIoT) catalogue its promise and its open problems \cite{essaid2025iiotbc}. Studies that combine blockchain with Zero Trust have followed \cite{liu2024dissecting,nie2025zeroTrust,salim2026ztblociot,mukta2026ztdelegation,jamwal2026survey,farhan2026ztfsf,potturi2026hybrid}.

\subsection{Problem Statement and Gap}
Reading this literature with an operator's eye, three things are missing. First, most blockchain-IoT identity schemes still authenticate once and then treat the session as trusted, which is precisely what Zero Trust is meant to remove \cite{wang2026medicaliot,james2024authz}; and the Zero Trust proposals for industrial settings generally assume one organization owns the policy decision point \cite{kumar2024industry40,zanasi2024flexible}. Second, a single trust threshold is applied to every request, although reading a temperature and rewriting a safety-controller setpoint carry very different consequences; and, as we show analytically and experimentally, a fixed threshold on a weighted score can be unreachable by an attacker who spoofs most of the signals. Third, nearly all of these designs implicitly place every decision on the ledger, which cannot meet the rates and latency of plant networks, and none of them says how a manufacturer's firmware recall, a legacy device that cannot hold a key, or a lost connection to the ledger should be handled. We therefore ask:

\emph{How can a consortium-operated blockchain identity registry be combined with continuous, consequence-aware authorization so that (i) revocation, including large-scale firmware recalls, takes effect without a token or list-refresh window and without any single party being able to impose it, (ii) the strength of verification scales with the criticality of the resource, and (iii) the enforcement path stays fast, available and auditable on industrial networks?}

\subsection{Contributions}
This paper makes the following contributions.
\begin{itemize}[leftmargin=1.1em,itemsep=0pt,topsep=1pt]
  \item \textbf{Architecture.} A zone-aware design (Section~\ref{sec:arch}) with three smart contracts (identity registry, Zero Trust access policy, quorum revocation governance), a tiered policy model that ties the required trust to the criticality of the target resource, an anomaly hard gate and a strict mode for safety-relevant writes, a protocol-gateway identity proxy for legacy devices, and batch revocation and quarantine for firmware recalls.
  \item \textbf{Hybrid enforcement.} A two-path decision model: a fast path in which a zone gateway decides locally against replicated ledger state and anchors a Merkle root of its decision log, and an audited path in which high-consequence requests are decided by the contract itself (Section~\ref{sec:hybrid}).
  \item \textbf{Analysis.} A threat model for plant environments and six proved properties, including a bound on how long a revoked device can still be accepted on the fast path and a result that exactly characterizes what the anomaly gate does and does not stop (Sections~\ref{sec:model}--\ref{sec:security}).
  \item \textbf{Open implementation and measurement.} Solidity 0.8.24 contracts executed on an EVM with measured gas for every operation, including recall and anchoring, plus a Python testbed using real Ed25519 signatures. Experiments report confidence intervals, a revocation-exposure comparison against token and CRL baselines, per-tier attack and false-denial rates, a gate sensitivity study, quorum-governance behaviour for 5--21 voters and recall capacity (Section~\ref{sec:eval}). A further evaluation replaces the synthetic behavioural signal with detector scores on real TON\_IoT traffic, including attack types withheld from training and an adaptive attacker (Section~\ref{sec:real_eval}), repeated on a SCADA dataset (Section~\ref{sec:wustl_eval}) and on an industrial dataset with labelled attack types (Section~\ref{sec:wustl21_eval}). The contracts are also deployed on a four-validator Hyperledger Besu QBFT network, first on one host and then with one validator in each of four Google Cloud regions, where we measure finality, revocation visibility, throughput, quorum voting, batch recall and validator failure (Section~\ref{sec:besu_eval}). We state plainly which results are compute-bound and which remain synthetic.
\end{itemize}
